\section{Discussion}\label{sec:discussion}

This paper evaluates one integrated system end to end: a physics-grounded wall closure connected
to a conditional generator.
Reading matching-height flow state, the closure predicts held-out traction better than its
equilibrium law ($0.540$ against $0.413$); supplying it improves source-excluded fluctuation
$R^2$ by $+0.061\,[0.057,0.064]$ under a pre-fixed rule. The unit was prospectively unscored,
not independent.

On the separating hill, the diffusion positive control passes and closure traction gains
$+0.062\,[+0.017,+0.106]$, not resolved from exact traction and better than wrong-time
traction. Roughly three effective samples make this indicative rather than precise.

One preregistered arm was adverse: flow matching fails its positive control and cannot adjudicate
the interface. We tested the obvious under-dispersion explanation prospectively with a stochastic
sampler on byte-identical weights; dispersion improved while the control failed \emph{more}
clearly. A model-family difference in how sparse wall information is used remains plausible, but
no mechanism is established. Ensemble dispersion is therefore reported as a diagnostic, never an
acceptance criterion (Supplementary Note~10).

Two properties make the cube number interpretable: its decisive
control is a real traction of the wrong instant, scoring worse than supplying nothing, so the
model is rewarded for \emph{this instant's} wall state; and the gain is distributional, the
closure arm best on every proper score at every ensemble size.

The practical reading is specific: expanding wall stress into an equilibrium-law velocity band
and clamping it fails, whereas the same quantity transmits through conditioning trained to read
it. The instantaneous field a generator needs is not the field an equilibrium law describes.

The limits are stated once, each marking where the evidence ends and the next study begins. The
composition is established offline and a priori, never inside a time-evolving solver.
Complete-volume absolute skill on the cube record stays negative for every
arm ($-0.028$ closure against $-0.088$ absence)---wall information moves the field towards the
target without yet making it accurate---while the re-test's near-wall region and the channel
record reach positive absolute skill. All three records carry few effective independent events,
so every interval is conditional on them---the price of difficult engineering records over
abundant homogeneous ones. The farther-from-wall gain, $+0.009\,[0.003,0.016]$, is
the honest measure of reach; the out-of-training probes are sensitivity checks, not causal
controls; and the fitted target is a first-cell viscous traction on the model raster, not the
native spectral-element gradient. The closure reads its matching heights at fixed multiples of
the local cell size, $2.5\Delta$ and $4.5\Delta$---the seam a wall-modelled solver carries---so
the multiples and the mesh setting $\Delta$ are parameters of the interface itself, and the
natural first task for a deployment study is charting the gain's response to coarser rasters and
different seams. The same standard then extends to the regimes most likely to strain the
interface: unresolved near-wall bands, spanwise three-dimensional near-wall structure, rough
walls and Reynolds numbers far above these records, where the wall's far-field footprint
grows.

The supportable engineering consequence is offline and specific: the implied streamwise viscous
wall-force error falls by almost two-fifths with the closure, while the equilibrium closure is
worse than absence on this endpoint despite a similar field-$R^2$ gain---a wall-load criterion
discriminates between closures that a volume-averaged score does not. Additional closure-side
geometry-transfer, separation, two-to-three-dimensional-transfer and diversity-scaling strengths
come from separate frozen companion models preserved by unchanged source hashes; the present
generator compositions do not revalidate those claims.

Wall-observation studies show that deployable pressure and shear signals can inform turbulent
fields \cite{guemes2021coarse,cuellar2024three,parikh2026conditional}, and field-generative
studies demonstrate wall-bounded generation, super-resolution and sparse reconstruction
\cite{gao2023bayesian,du2024confild,oommen2026learning}. This work complements them by taking
the wall quantity from a closure rather than an observation, excluding supplied and read
cells from every score, and testing against an in-distribution wrong-information control. It
establishes no superiority over prior architectures: no head-to-head comparison was run.

What a reader can take from this is separable from our own closure. The conditioning slot is a
harness, not a result: a wall model supplying a signed two-component wall-on-fluid traction enters as a
surface field on the wall faces alone, the generator never sees which model occupies it, and a
sublayer projection makes the generated wall gradient equal the supplied traction identically, so
two wall models compete with no generator-side advantage (Fig.~\ref{fig:slot}). The
fidelity--gain response measured through it, monotone from $+0.104$ to $+0.153$ across traction
fidelities $0.082$--$0.797$, states how accurate a wall model must become before a generative
reconstruction benefits; its limit is equally usable, since a data-only regressor of better
scalar fidelity gains more than the ladder predicts, so error \emph{structure} matters beyond any
scalar. We propose the oracle positive control as the minimum adequacy standard such a study needs
before reading a conditioning contrast, and the grouped physical-time protocol, with the $0.839$ inflation a naive
flattened split produces, applies to anyone evaluating this widely used public corpus. All are
released with the code and data.

The three records show that generative models do not bypass the wall problem: they demonstrate
the framework on pinned, attached and separating cases, not a claim to have covered wall
modelling. What remains is runtime composition: a solver-coupled calculation supplying a momentum
boundary condition as the flow evolves. Transmission holds in both generative families on the
pinned cube and attached channel records, and in one of the two on the separating record,
the other an adverse preregistered arm for which no mechanism is claimed. Within those bounds a
physics-grounded wall closure can supply information a generative model demonstrably propagates
through the field.
